\textbf{Coordinación de detección y liberación con hidrógeno (INC-21).}\label{sec:sd4-inc21}
Las redes F04/F05 conservan ocho puntos y 14 m por banco, pero se retienen sus dos VEP-A00-1P-UL seleccionados en INC-30 y confirmadores pendientes, sin autorización de energización hasta completar la clasificación y el expediente. La clasificación debe abarcar emisión normal, extremo conservador de recarga de 60 A por UPS no habilitado conforme a GAS-28, pérdida de extracción y emisión posterior al aislamiento: la renovación de 900 m$^3$/h no permite declarar por sí sola todo el banco no clasificado. El expediente deberá identificar zona y extensión, grupo de gas compatible con hidrógeno, temperatura y condiciones especiales del certificado; cubrirá cámara aspirante, entrada, escape, filtros, tubos, uniones, sellos, prensaestopas, confirmador independiente y equipos eléctricos de la zona. Un certificado del gabinete no se extiende automáticamente a sus componentes ni a una mezcla que entra por aspiración. Las seis estaciones manuales y abortos se sitúan fuera del volumen de riesgo, en el lado seguro de cada salida; el contacto de confirmación de banco deberá entrar al monitor supervisado Fike sin introducir una fuente de ignición. La ficha VLF-500-00-UL queda como antecedente de las cuatro zonas ordinarias, sustituidas por VEP en INC-36. El conjunto de bancos no se reemplaza por un modelo con designación comercial Ex sin certificado y manual verificables. La selección apta pendiente mantiene abierto el diseño integral de INC-4.

Para EXT-4 se mantiene extracción segregada de bancos durante alarma, aislamiento y posaislamiento. No se cierra una compuerta de banco ni se autoriza inundación de éste a partir de una lectura instantánea de H$_2$ baja. Con $V=36$ m$^3$, $Q=900$ m$^3$/h y mezcla ideal, la concentración de agente sin aporte adicional evoluciona como $C(t)=C_0e^{-Qt/V}$. A diez minutos, $e^{-900(10/60)/36}=0,01550$: una concentración inicial de 8,0\,\% caería a 0,1240\,\%. Esa renovación elimina el fundamento de una inundación única con retención de diez minutos. No constituye simulación CFD ni demuestra concentraciones locales. Se descarta cerrar ventilación mientras se mantiene la emisión conservadora de 7,5978 m$^3$/h. Si la emisión residual fuera constante y se sellara el banco desde 0,4\,\% con límite propio de 1,0\,\%, el máximo compatible con diez minutos sería $(0,010-0,004)36/(10/60)=1,296$ m$^3$/h; representa una reducción de al menos 82,94\,\% frente al escenario conservador. Ese límite no acredita la emisión de los bloques CSB ni autoriza sellado: se requiere su curva residual, medición distribuida y estudio gas/fuego. La confirmación de fuego ordena bloqueo de recarga y aislamiento mediante dispositivo DC apto, evacuación y aviso; no depende de WAN ni se considera cumplida por abrir sólo la entrada del rectificador.

\textbf{Compatibilidad y reserva física actualizada EXT-30.}\label{sec:sd4-ext21}
EXT-30 sustituye las referencias físicas DOT/EN previas por una cohorte uniforme EN del manual Fike 06-542 UL/FM: 83 L 70-256, dos 26 L 70-278 y tres 16 L 70-350, más seis reemplazos. Los llenados son 55/24/24/11/11/15 kg por cohorte; tara más agente suma 592 kg entre doce. Se seleccionan DFIA/RCM e interruptores correspondientes al mismo conjunto. Las dimensiones de fichas DOT anteriores no definen compra EN. Se mantienen dos reservas exteriores de gabinete 3 por 0,8 por 1,8 m con frente libre 1,20 m; conexiones, soporte, control térmico, listado vigente, hidráulica y seguridad gas/fuego permanecen por verificar antes de habilitar. Las masas y dimensiones actuales se desarrollan en EXT-30.

\textbf{Protección secundaria portátil PORT-21.}\label{sec:sd4-port21}
Synaptix selecciona veintinueve extintores instalados y cinco de reserva local, total treinta y cuatro; PORT-24 fija posiciones y cobertura vigente. El Amerex B441 utiliza polvo químico seco ABC, 10 lb (4,536 kg), potencial publicado 4A:80B:C y presión nominal 195 psig; el Amerex 398 emplea Halotron I, 15,5 lb (7,031 kg), 2A:10B:C; el Amerex 322 emplea CO$_2$, 5 lb (2,268 kg), 5B:C. Las clases y potenciales provienen del catálogo del fabricante; su listado UL/ULC no se transforma en certificación chilena del lote. Cada modelo conservará rotulación, boquilla/manguera y soporte originales, manual en español y certificación exigible al habilitar \parencite[pp.~9, 11--12]{l21-amerex}.

\begin{longtable}{|p{0.12\linewidth}|p{0.19\linewidth}|p{0.54\linewidth}|}
\hline \textbf{Identidad} & \textbf{Modelo / cantidad} & \textbf{Ubicación prescrita y amenaza} \\\hline
\endfirsthead
\hline \textbf{Identidad} & \textbf{Modelo / cantidad} & \textbf{Ubicación prescrita y amenaza} \\\hline
\endhead
P01--P06 & B441 / 6 & Uno por zona F01--F06, en lado seguro de su salida. F04/F05: exterior norte próximo a puerta, sin ingresar al banco para obtenerlo; F02/F03: exterior oeste/este junto a salida. F01: junto a puerta sur hacia circulación; F06: junto a acceso sur. Cobertura A de materiales/cableado, B de líquidos y C de equipo energizado. \\\hline
P07 & 398 / 1 & Interior F01, junto a salida sur, para intervención incipiente autorizada en TI; conserva el B441 para cobertura A mínima. No se sitúa en UPS/bancos: el fabricante exige volumen de almacenamiento mínimo 2.170 pies$^3$ = 61,45 m$^3$, que F01 (82,5 m$^3$) satisface y los recintos pequeños no satisfacen. \\\hline
P08/P09 & 322 / 2 & Lado seguro exterior junto a cada salida de UPS, separados del B441 y señalizados; opción de intervención sin residuo en equipo energizado. No se ingresa a un recinto inundado con CO$_2$; éste no sustituye cobertura A ni autoriza intervención frente a gas inflamable. \\\hline
P10--P29 & B441 / 20 & Estaciones exteriores, coordenadas y recorridos calculados en PORT-24 y Figura~\ref{fig:recinto-exterior-coordinado}. Abarcan energía, combustible, bancadas, climatización, UTA, agua, recipientes, trabajo y stock. \\\hline
R-PQS/R-HAL/R-CO2 & B441 / 3; 398 / 1; 322 / 1 & Cinco reservas cargadas en almacén local separado, protegidas y con certificación vigente; no se cuentan como posiciones activas. Reemplazo del mismo tipo antes de retirar un equipo. \\\hline
\end{longtable}

La cobertura y recorridos actuales se desarrollan en PORT-24: no se conserva como acreditación la suma nominal de lados del plano anterior de 30 por 33 m. El módulo interior y cada sector exterior conservan protección propia, de acuerdo con potencial y distancia del DS 594 \parencite[arts.~45--49]{l21-ds594}.

Los B441 ocupan aproximadamente 0,521$\times$0,222$\times$0,127 m; el 398, 0,556$\times$0,235$\times$0,152 m; el 322, 0,451$\times$0,210$\times$0,133 m. Se reservan nichos de 0,35$\times$0,25 m de planta y 0,70 m de alto por unidad, fuera de barridos y sin reducir paso útil bajo 1,20 m; se instala la base a máximo 1,30 m del suelo y señal visible. Los equipos exteriores requieren gabinete contra intemperie/corrosión que permita retiro inmediato sin llave. Las reservas necesitan cinco nichos iguales en el almacén; su consumo eléctrico es nulo; las veintiséis unidades B441 activas requieren tres reservas por redondeo superior del diez por ciento, además de una por cada otro tipo \parencites[pp.~9, 11--12]{l21-amerex}[arts.~47/49]{l21-ds594}.

Synaptix asigna al responsable local de operación una inspección documentada mensual y después de cada incidente: acceso, sello/pasador, manguera/boquilla, daños/corrosión, carga por pesada cuando corresponda, presión, etiqueta y vencimiento. Un servicio técnico certificado conforme al DS 44, NCh3268 y NCh2056 ejecutará mantenimiento al menos anual, recarga después de uso y pruebas hidrostáticas según norma/manual; se verifica su certificado y alcance antes de contratar, sin atribuir un contrato ya suscrito. El registro identifica serial, posición P/R, modelo, agente, masa, potencial local acreditado, certificado del producto y prestador, fecha, actuación, técnico y próximo vencimiento. El responsable local instala reemplazo equivalente antes de retirada, sin consumir protección de otro sector; las reservas se reponen inmediatamente y cualquier indisponibilidad produce aviso NOC/CLIENTE y vigilancia compensatoria hasta restablecimiento. La habilitación exige inventario, certificados vigentes, etiquetas en español, manual, plano medido y acta de accesibilidad. La capacitación de todo personal local en alarma, evacuación y uso autorizado precede la puesta en servicio; los portátiles no sustituyen extinción automática ni autorizan a personal sin entrenamiento a combatir incendio de bancos \parencites[arts.~23--30/32--33]{l21-ds44}[arts.~48/51]{l21-ds594}.

\textbf{Cobertura calculada.} PORT-24 reemplaza el mapa preliminar por rutas alrededor de huellas y corredores propios dimensionados, incluyendo las UTA superiores y los puntos de trabajo/stock. La certificación chilena, capacitación y servicio se entregan en el hito de habilitación; no se presentan como ejecutados.
